\section{Ordre des octets : \emph{big endian} et \emph{little endian}}
\label{secEndianness}

Dès qu'une information occupe \textbf{plusieurs octets} (un entier
16~ou~32~bits, une position cartésienne codée en réel), la conversion ne
se limite plus au découpage bit à bit vu dans la section précédente : il
faut aussi savoir dans quel \textbf{ordre} ranger ces octets en mémoire.

\begin{UPSTIinfor}{\emph{Big endian} / \emph{little endian}}
Pour une donnée multi-octets, deux conventions existent :
\begin{itemize}
    \item \textbf{Big endian} : l'octet de poids \textbf{fort} est stocké
    en premier (à l'adresse la plus basse).
    \item \textbf{Little endian} : l'octet de poids \textbf{faible} est
    stocké en premier.
\end{itemize}
Deux équipements qui échangent une donnée multi-octets doivent
s'accorder sur cette convention, sous peine de lire une valeur
totalement erronée bien que chaque octet, pris isolément, ait été reçu
correctement.
\end{UPSTIinfor}

\textbf{Exemple.} La valeur entière \texttt{0x1234} (4660 en décimal),
codée sur 2~octets, est rangée ainsi :

\begin{center}
\begin{tabular}{|l|c|c|}
\hline
 & \textbf{Octet à l'adresse $n$} & \textbf{Octet à l'adresse $n+1$} \\
\hline
Big endian & \texttt{0x12} & \texttt{0x34} \\
\hline
Little endian & \texttt{0x34} & \texttt{0x12} \\
\hline
\end{tabular}
\end{center}

Un équipement qui interprèterait en \emph{big endian} une donnée envoyée
en \emph{little endian} lirait \texttt{0x3412} (13330) au lieu de
\texttt{0x1234} (4660).

\begin{UPSTIinfor}{Application sur la cellule : configuration du serveur uniVALplc}
Le contrôleur CS9 du robot encode les données échangées via le serveur
\texttt{uniVALplc} selon l'ordre d'octets attendu par le client. Comme
l'automate client de la cellule est un \textbf{M340 de Schneider}
(processeur little endian), l'option \textbf{\emph{Little Endian}} a été
sélectionnée lors de l'installation du serveur \texttt{uniVALplc} sur la
baie CS9 : sans cette configuration, tous les champs numériques
multi-octets de \texttt{T\_StaeubliRobot} (vitesses, positions
cartésiennes, compteurs...) seraient mal interprétés par le programme
Control Expert, même si les blocs \texttt{VAL\_ReadAxesGroup} et
\texttt{VAL\_WriteAxesGroup} sont correctement appelés.
\end{UPSTIinfor}

\begin{UPSTIactivite}[4][Lecture d'une valeur multi-octets][][][TD]
\UPSTIquestion{Le champ \texttt{Status.Heartbeat} de
\texttt{T\_StaeubliRobot} est de type \texttt{UINT}, codé sur 2~octets.
Le coupleur \texttt{NOC\_ETHIP\_ROBOT} reçoit, dans l'ordre, les deux
octets \texttt{0x64} puis \texttt{0x00}. Sachant que le serveur
\texttt{uniVALplc} a été configuré en \emph{little endian}, calculer la
valeur numérique (en décimal) de \texttt{Status.Heartbeat}.}
\UPSTIquadrillage{2.5}

\begin{UPSTIcorrectionP}{Corrigé}
En \emph{little endian}, le premier octet reçu est l'octet de poids
\textbf{faible} : \texttt{0x64} (poids faible) et \texttt{0x00} (poids
fort) donnent \texttt{0x0064}, soit
$6\times16+4 = \mathbf{100}$.
\end{UPSTIcorrectionP}

\UPSTIquestion{Un collègue, en reprenant la configuration du serveur sur
une autre cellule, laisse par erreur l'option \emph{Big Endian}. Quelle
valeur le programme Control Expert lirait-il alors pour
\texttt{Status.Heartbeat}, à partir des mêmes octets \texttt{0x64} puis
\texttt{0x00} ? Ce défaut de configuration serait-il facile à détecter au
premier démarrage de la cellule ? Justifier.}
\UPSTIquadrillage{3.5}

\begin{UPSTIcorrectionP}{Corrigé}
En \emph{big endian}, le premier octet reçu est interprété comme l'octet de
poids \textbf{fort} : le programme lirait \texttt{0x6400}, soit
$100\times256 = \mathbf{25\,600}$.

Ce défaut est \textbf{difficile à détecter} au premier démarrage : chaque
octet est correctement reçu, les bits de santé du coupleur
\texttt{NOC\_ETHIP\_ROBOT} restent à \texttt{TRUE} et les blocs
\texttt{VAL\_ReadAxesGroup}/\texttt{VAL\_WriteAxesGroup} sont appelés
normalement : aucune erreur de communication n'est signalée. Seules des
valeurs numériques \textbf{incohérentes} trahissent le problème (par
exemple une position cartésienne lue sans rapport avec celle affichée sur
le pendant du robot, une vitesse absurde) -- encore faut-il penser à les
comparer. Un champ comme \texttt{Heartbeat}, dont la valeur absolue n'a
pas de signification pour l'opérateur, ne permet pas à lui seul de s'en
apercevoir.
\end{UPSTIcorrectionP}
\end{UPSTIactivite}
